\documentclass[a4paper,12pt]{article}
\usepackage[T2A]{fontenc}
\usepackage[utf8]{inputenc}
\usepackage[russian]{babel}
\usepackage{amsmath,amssymb,mathtools}
\usepackage{helvet}
\renewcommand{\familydefault}{\sfdefault}
\usepackage{icomma}
\usepackage[margin=20mm,top=22mm,bottom=22mm]{geometry}
\usepackage{microtype}
\usepackage{tikz}
\usetikzlibrary{arrows.meta,positioning,shapes.geometric,calc}
\usepackage{tabularx,booktabs,array}
\usepackage{enumitem}
\usepackage{parskip}
\usepackage{fancyhdr}
\usepackage{setspace}
\usepackage{xcolor}
\usepackage{hyperref}
\definecolor{accent}{HTML}{008080}
\definecolor{darkaccent}{HTML}{00595B}
\definecolor{textgray}{HTML}{333333}
\definecolor{muted}{HTML}{687778}
\hypersetup{colorlinks=true,linkcolor=darkaccent,bookmarksnumbered=true,pdfauthor={Лёвин Артём Александрович},pdftitle={Чек-лист. Теория вероятностей и комбинаторика}}
\onehalfspacing
\setlength{\parindent}{0pt}
\setlength{\parskip}{6pt}
\setlist[itemize]{leftmargin=1.6em,itemsep=2pt,topsep=3pt}
\setlist[enumerate]{leftmargin=1.65em,itemsep=5pt,topsep=4pt}
\newcommand{\key}[1]{\textcolor{darkaccent}{\textbf{#1}}}
\newcommand{\sample}[1]{\par\medskip\textcolor{darkaccent}{\textbf{Пример.}}\quad #1\par}
\pagestyle{fancy}
\fancyhf{}
\fancyhead[L]{\small\color{muted}Теория вероятностей и комбинаторика}
\fancyhead[R]{\small\color{muted}Екатерина Гнедкова}
\fancyfoot[C]{\small\color{muted}\thepage}
\renewcommand{\headrulewidth}{0.25pt}
\renewcommand{\footrulewidth}{0pt}
\setlength{\headheight}{16pt}
\newcommand{\sect}[1]{\section{#1}\vspace{-3pt}}
\newcommand{\subsect}[1]{\subsection*{\textcolor{darkaccent}{#1}}\vspace{-8pt}}
\begin{document}
\thispagestyle{empty}
\begin{center}
\vspace*{30mm}
{\small\color{darkaccent}\MakeUppercase{Математика}}\par
\vspace{17mm}
{\LARGE Чек-лист}\par
\vspace{10mm}
{\Large\color{accent}Теория вероятностей\\[2mm]и комбинаторика}\par
\vspace{16mm}
{\normalsize Подготовка к ЕГЭ по профильной математике}\par
\vspace{33mm}
{\normalsize Екатерина Гнедкова}\par
\vspace{8mm}
{\normalsize 4 ноября 2026 г.}\par
\vspace{20mm}
{\small Лёвин Артём Александрович\\ эксклюзивно для Екатерины Гнедковой}\par
\vfill
\begin{tikzpicture}
  \draw[accent,line width=0.8pt] (-5,0) .. controls (-3.7,0.6) and (-1.7,-0.6) .. (0,0)
  .. controls (1.8,0.6) and (3.4,-0.6) .. (5,0);
  \fill[accent] (0,0) circle (1.0pt);
\end{tikzpicture}
\vspace*{13mm}
\end{center}
\clearpage

\sect{Классическая вероятность}
\key{Элементарный исход} --- один возможный результат случайного опыта. При условии, что все исходы равновозможны, вероятность события $A$ вычисляется по формуле
\[
\boxed{P(A)=\frac{m}{n}},
\qquad n\text{ --- число всех исходов},\quad m\text{ --- число благоприятных.}
\]
Всегда выполняется $0\le P(A)\le1$. Сначала определяйте, \emph{что именно считается одним исходом}: порядок людей, последовательность бросков или состав выбранной группы.

\sample{На телефоне случайно нажата одна из десяти цифр от $0$ до $9$. Чётных цифр пять: $0,2,4,6,8$. Поэтому $P=5/10=1/2$. Нуль тоже чётный.}

\sect{Метод позиций и факториал}
\key{Правило произведения.} Если первый шаг выполняется $a$ способами и после каждого из них второй --- $b$ способами, то последовательность из двух шагов можно выполнить $a\cdot b$ способами. Для большего числа шагов рассуждаем аналогично.

\sample{Три человека выстраиваются в очередь. На первое место есть три претендента, на второе --- два, на третье --- один. Получаем $3\cdot2\cdot1=3!=6$ очередей.}

Факториал натурального числа:
\[
n!=n(n-1)\cdots2\cdot1,\qquad 0!=1.
\]
\key{Когда порядок важен.} Очереди $ABC$ и $BAC$ считаются разными. Результаты трёх бросков $(1,2,3)$ и $(3,2,1)$ тоже различны.

\key{Когда порядок не важен.} Группы $\{A,B\}$ и $\{B,A\}$ совпадают. Если выбираем $k$ различных участников из $n$ без возвращения, число групп равно
\[
\boxed{C_n^k=\frac{n!}{k!(n-k)!}=\frac{n(n-1)\cdots(n-k+1)}{k!}}.
\]
Деление на $k!$ убирает повторный учёт одной группы во всех перестановках её участников.

\sample{Из пяти туристов выбирают двоих. Всего $C_5^2=5\cdot4/2!=10$ пар. Если турист $A$ должен войти в пару, второго выбирают четырьмя способами. Значит, $P=4/10=2/5$.}

\clearpage
\sect{Выбор с обязательным участником}
Из $n$ человек случайно выбирают группу из $k$ человек. Если конкретный человек $A$ должен оказаться в группе, он уже зафиксирован, а остальных $k-1$ выбирают из $n-1$ человек:
\[
N=C_n^k,\qquad M=C_{n-1}^{k-1},\qquad\boxed{P=\frac{M}{N}=\frac{k}{n}}.
\]
\sample{Из $20$ туристов трое уходят в село. Вероятность того, что с ними пойдёт турист $A$:
\[
P=\frac{C_{19}^{2}}{C_{20}^{3}}=
\frac{19\cdot18/2!}{20\cdot19\cdot18/3!}=\boxed{\frac{3}{20}}.
\]
Делим на $2!$ в числителе: турист $A$ уже включён, и порядок имеет значение только при выборе \emph{двух его спутников}. Формула $k/n$ верна для равновероятного выбора группы без дополнительных ограничений.}

\sect{Условие на относительный порядок и соседство}
\subsect{Кто окажется раньше?}
В случайной очереди из трёх человек $A,B,C$ вероятность того, что $B$ будет позже $A$, равна $1/2$. Из каждой очереди, где $A$ перед $B$, перестановкой этих двоих получаем единственную очередь, где $B$ перед $A$. Следовательно, случаев поровну.

\subsect{Два человека сидят рядом}
Если из $r$ различных людей два заданных человека обязаны стоять рядом, рассматриваем их как \emph{один блок}. Блок и остальные $r-2$ человека можно переставить $(r-1)!$ способами, а внутри блока порядок может быть любым из $2!$:
\[
P(\text{рядом})=\frac{2!(r-1)!}{r!}=\boxed{\frac{2}{r}},\quad r\ge2.
\]
\sample{Для трёх человек $P=2/3$, для четырёх --- $P=1/2$.}

\clearpage
\sect{Вероятность через дополнение}
Иногда выгоднее вычесть неподходящие случаи из общего числа. Всегда верно:
\[
P(A)=1-P(\overline A),
\]
где $\overline A$ --- событие «$A$ не произошло».

\subsect{Диагональ многоугольника}
Из $r$ вершин выпуклого $r$-угольника выбирают две. Всего пар $C_r^2$. Пара образует сторону, если выбраны соседние вершины; таких пар ровно $r$. Остальные пары задают диагонали:
\[
M=C_r^2-r,\qquad
P(\text{диагональ})=\frac{C_r^2-r}{C_r^2}=
\boxed{\frac{r-3}{r-1}}.
\]
\sample{У четырёхугольника $C_4^2=6$ пар, из которых две --- диагонали: $P=2/6=1/3$. У двадцатиугольника $P=17/19$.}

\sect{Когда возможны повторения}
Если два человека независимо выбирают по одной цифре из $\{1,2,3\}$ и каждый может выбрать любую из трёх цифр, всего $3\cdot3=9$ \emph{упорядоченных} результатов. Совпадение разрешено.

\sample{Вероятность разных цифр: первый выбирает любую из трёх, второй --- одну из двух оставшихся. Поэтому $P=(3\cdot2)/(3\cdot3)=2/3$.}

\key{Важное различие:} при независимых выборах с возможным повторением число результатов $n^k$; при выборе $k$ разных предметов без возвращения и без учёта порядка --- $C_n^k$.

\sect{Два броска игрального кубика}
Для честного кубика исход каждого броска --- число от $1$ до $6$. При двух бросках учитываем порядок: $(2,5)$ и $(5,2)$ --- разные результаты. Всего $6^2=36$ равновозможных исходов.
\begin{itemize}
\item Первый бросок равен $6$: $6$ подходящих исходов, $P=6/36=1/6$.
\item Числа совпали: $(1,1),\ldots,(6,6)$; $P=6/36=1/6$.
\item Сумма чётная: оба числа чётные или оба нечётные; $M=3\cdot3+3\cdot3=18$, $P=1/2$.
\end{itemize}

\clearpage
\sect{Таблица исходов для двух бросков}
Если спрашивают о сумме очков, удобна таблица: строки соответствуют первому броску, столбцы --- второму. В ячейке записана сумма.

\begin{center}
\renewcommand{\arraystretch}{1.45}
\setlength{\tabcolsep}{11pt}
\begin{tabular}{c@{\hspace{8pt}}cccccc}
\toprule
& \multicolumn{6}{c}{\small Второй бросок}\\
\cmidrule(l){2-7}
\small Первый & $1$ & $2$ & $3$ & $4$ & $5$ & $6$ \\
\midrule
$1$ & $2$ & $3$ & $4$ & $5$ & $6$ & \textcolor{darkaccent}{\textbf{7}}\\
$2$ & $3$ & $4$ & $5$ & $6$ & \textcolor{darkaccent}{\textbf{7}} & $8$\\
$3$ & $4$ & $5$ & $6$ & \textcolor{darkaccent}{\textbf{7}} & $8$ & $9$\\
$4$ & $5$ & $6$ & \textcolor{darkaccent}{\textbf{7}} & $8$ & $9$ & $10$\\
$5$ & $6$ & \textcolor{darkaccent}{\textbf{7}} & $8$ & $9$ & $10$ & $11$\\
$6$ & \textcolor{darkaccent}{\textbf{7}} & $8$ & $9$ & $10$ & $11$ & $12$\\
\bottomrule
\end{tabular}
\end{center}

\key{Самая вероятная сумма --- $7$.} Она встречается в шести клетках из $36$:
\[
P(\text{сумма }7)=\frac{6}{36}=\frac16.
\]
Таблица помогает увидеть отдельные случаи и проверить подсчёт. При большом числе исходов обычно удобнее метод позиций.

\sect{Три броска: похожие, но разные события}
Всего упорядоченных результатов $6^3=216$.

\renewcommand{\arraystretch}{1.35}
\begin{center}
\small
\begin{tabularx}{\linewidth}{@{}X>{\raggedleft\arraybackslash}p{34mm}>{\raggedleft\arraybackslash}p{22mm}@{}}
\toprule
Событие & Благоприятных & Вероятность\\
\midrule
Все три числа одинаковы & $6\cdot1\cdot1=6$ & $\frac{1}{36}$\\
Все три числа различны & $6\cdot5\cdot4=120$ & $\frac{5}{9}$\\
Только первые два одинаковы, третье отличается & $6\cdot1\cdot5=30$ & $\frac{5}{36}$\\
Ровно два из трёх одинаковы, на любых позициях & $3\cdot6\cdot5=90$ & $\frac{5}{12}$\\
\bottomrule
\end{tabularx}
\end{center}

В последней строке множитель $3$ --- число способов выбрать позицию единственного отличающегося броска: первый, второй или третий. Можно проверить дополнением:
\[
216-6-120=90.
\]
\key{Не путайте «ровно два» и «хотя бы два».} Во втором случае подходят также три одинаковых числа: $90+6=96$ исходов, то есть вероятность $96/216=4/9$.

\clearpage
\begin{center}
{\large\bfseries\color{darkaccent}Алгоритм: как считать равновозможные исходы}
\end{center}
\vspace{8mm}
\begin{center}
\begin{tikzpicture}[
 >=Stealth,
 line/.style={-Stealth,draw=darkaccent,line width=0.85pt},
 start/.style={ellipse,draw=accent,rounded corners=1pt,align=center,minimum width=45mm,minimum height=10mm,inner sep=4pt},
 action/.style={rectangle,rounded corners=2pt,draw=accent,align=center,text width=84mm,minimum height=13mm,inner sep=5pt},
 branch/.style={rectangle,rounded corners=2pt,draw=accent,align=center,text width=55mm,minimum height=23mm,inner sep=6pt},
 decide/.style={diamond,aspect=2.1,draw=accent,align=center,text width=39mm,inner sep=2pt,minimum height=15mm}
]
\node[start] (begin) at (0,0) {Начало};
\node[action] (identify) at (0,-2) {Определите, что считается одним исходом.\\Убедитесь в равновозможности исходов.};
\node[decide] (decision) at (0,-4.45) {Важен ли\\ порядок?};
\node[branch] (yes) at (-4.05,-7.4) {Да. Метод позиций.\\При независимых выборах с повторением: $n^k$.};
\node[branch] (no) at (4.05,-7.4) {Нет. Для $k$ различных элементов из $n$ без возвращения: $C_n^k$.};
\node[action] (count) at (0,-10.25) {Таким же способом посчитайте\\благоприятные исходы $M$ и все исходы $N$.};
\node[start] (end) at (0,-12.3) {$P=\dfrac{M}{N}$};
\draw[line] (begin)--(identify);
\draw[line] (identify)--(decision);
\draw[line] (decision.west)-| node[pos=0.2,above]{\small да} (yes.north);
\draw[line] (decision.east)-| node[pos=0.2,above]{\small нет} (no.north);
\coordinate (join) at (0,-9.15);
\draw[draw=darkaccent,line width=0.85pt] (yes.south) |- (join);
\draw[draw=darkaccent,line width=0.85pt] (no.south) |- (join);
\draw[line] (join)--(count.north);
\draw[line] (count)--(end);
\end{tikzpicture}
\end{center}
\vfill
\begin{center}\small\color{muted}Если исходы не равновозможны, формулу $M/N$ применять напрямую нельзя.\end{center}
\clearpage

\section*{Тренировочный блок. Путь А}
\addcontentsline{toc}{section}{Тренировочный блок}
Во всех задачах исходы считаются равновозможными; там, где указано, выбор делается случайно и независимо. Покажите подсчёт общего числа исходов и благоприятных исходов.

\subsect{Средний уровень}
\begin{enumerate}
\item На клавиатуре телефона случайно нажимают одну цифру от $0$ до $9$. Какова вероятность нажать цифру, кратную трём? Нуль учитывайте.
\item Четыре человека, среди них Анна и Борис, случайно становятся в очередь. Найдите вероятность того, что Борис стоит позже Анны.
\item Из $10$ туристов случайно выбирают троих для похода. Какова вероятность того, что среди них окажется конкретный турист $A$?
\end{enumerate}

\subsect{Уровень выше среднего}
\begin{enumerate}[resume]
\item Из вершин выпуклого $12$-угольника наугад выбирают две различные. Какова вероятность того, что соединяющий их отрезок --- диагональ?
\item Пятеро различных людей садятся в ряд на пять стульев. Какова вероятность того, что Анна и Борис окажутся на соседних стульях?
\item Два человека независимо выбирают по одной цифре из набора $\{1,2,3,4\}$. Повторения разрешены. Найдите вероятность выбора различных цифр.
\item Честный игральный кубик бросают дважды. Какова вероятность того, что в первом броске выпадет $5$, а во втором --- чётное число?
\item Игральный кубик бросают дважды. Найдите вероятность того, что сумма очков равна $9$.
\end{enumerate}

\subsect{Сложный уровень}
\begin{enumerate}[resume]
\item Игральный кубик бросают трижды. Какова вероятность того, что все три выпавших числа различны?
\item Игральный кубик бросают трижды. Найдите вероятность того, что \emph{ровно два} из трёх выпавших чисел одинаковы.
\end{enumerate}

\clearpage
\section*{Ответы к тренировочному блоку}
\addcontentsline{toc}{section}{Ответы к тренировочному блоку}
\begin{center}
\renewcommand{\arraystretch}{1.7}
\begin{tabular}{@{}c@{\hspace{13pt}}c@{\hspace{20pt}}l@{}}
\toprule
№ & Ответ & Проверка подсчёта\\
\midrule
1 & $\displaystyle\frac25$ & $\displaystyle\frac4{10}$\\
2 & $\displaystyle\frac12$ & по симметрии порядка\\
3 & $\displaystyle\frac3{10}$ & $\displaystyle\frac{C_9^2}{C_{10}^3}$\\
4 & $\displaystyle\frac9{11}$ & $\displaystyle\frac{C_{12}^2-12}{C_{12}^2}$\\
5 & $\displaystyle\frac25$ & $\displaystyle\frac{2!\cdot4!}{5!}$\\
6 & $\displaystyle\frac34$ & $\displaystyle\frac{4\cdot3}{4\cdot4}$\\
7 & $\displaystyle\frac1{12}$ & $\displaystyle\frac{1\cdot3}{6\cdot6}$\\
8 & $\displaystyle\frac19$ & $\displaystyle\frac4{36}$\\
9 & $\displaystyle\frac59$ & $\displaystyle\frac{6\cdot5\cdot4}{6^3}$\\
10 & $\displaystyle\frac5{12}$ & $\displaystyle\frac{3\cdot6\cdot5}{6^3}$\\
\bottomrule
\end{tabular}
\end{center}

\vspace{10mm}
\subsect{Итоговая памятка}
\begin{itemize}
\item \key{Равновозможные исходы:} вероятность --- отношение подходящих исходов ко всем исходам.
\item \key{Очередь и броски:} важен порядок; используйте метод позиций.
\item \key{Группа без порядка:} используйте сочетания $C_n^k$.
\item \key{Фиксированный участник:} выбирайте остальных из оставшихся людей.
\item \key{Соседство:} объединяйте двух людей в блок и учитывайте внутренние перестановки.
\item \key{«Ровно два»:} исключайте случай, когда совпали все три.
\item \key{Сумма на двух кубиках:} таблица $6\times6$ помогает увидеть все варианты.
\end{itemize}
\vfill
\begin{center}\small\color{muted}Лёвин Артём Александрович эксклюзивно для Екатерины Гнедковой\end{center}
\end{document}
